\chapter{Limites et perspectives}
\label{chap:limites}

\section{Limites actuelles}
L'honnêteté scientifique impose d'identifier clairement les limites du travail réalisé.
\begin{itemize}[leftmargin=1.5em]
  \item \textbf{Validation sur données de référence.} Les performances sont établies sur des
        jeux de données publics (CWRU, CMAPSS, etc.) et des signaux de démonstration. Une
        validation sur des \emph{données de terrain réelles}, acquises en continu sur des
        machines en exploitation, reste à mener.
  \item \textbf{Transférabilité du pronostic.} Le modèle de RUL est entraîné sur des
        turbomoteurs simulés (CMAPSS). Sa transposition à un équipement client suppose de
        disposer de l'historique de dégradation \emph{de cet équipement} --- donnée que peu
        d'usines possèdent. Le pronostic doit donc être présenté comme une estimation, à
        recalibrer par machine.
  \item \textbf{Signaux de simulation.} En mode démonstration, les formes d'onde affichées
        sont reconstruites à partir des features ; l'ingestion temps réel de signaux bruts
        haute fréquence en continu (50\,kHz sur des centaines de machines) relève de la
        couche edge, prête mais non éprouvée à cette échelle.
  \item \textbf{Une tâche principale par dataset.} Chaque dataset est exploité pour sa tâche
        dominante ; une exploitation multi-tâche conjointe (diagnostic + pronostic sur le
        même équipement) reste à développer.
  \item \textbf{Connectivité éprouvée en laboratoire.} OPC-UA a été validé avec un serveur de
        démonstration, non avec un automate réel ; la bascule TimescaleDB est outillée mais
        non déployée en production.
  \item \textbf{Sécurité.} L'authentification JWT et la gestion de rôles constituent une
        base, mais un déploiement industriel exigerait un RBAC fin, l'intégration SSO/LDAP et
        un durcissement conforme à \textbf{IEC~62443}.
\end{itemize}

\section{Perspectives d'évolution}
Les perspectives s'organisent en trois horizons.

\subsection{Court terme --- validation et robustesse}
\begin{itemize}[leftmargin=1.5em]
  \item Campagne de validation sur banc d'essai instrumenté réel et collecte de données
        \emph{run-to-failure} pour recalibrer le pronostic.
  \item Mesure du taux de fausses alarmes en conditions réelles, via la boucle de validation
        expert déjà en place.
  \item Déploiement effectif sur TimescaleDB pour un parc de plusieurs centaines de machines.
\end{itemize}

\subsection{Moyen terme --- différenciation technologique}
\begin{itemize}[leftmargin=1.5em]
  \item \textbf{Modèles physics-informed} : hybrider les modèles de données avec des modèles
        physiques de dégradation (loi de Paris pour la fissuration, modèles de fatigue de
        roulement) pour améliorer la transférabilité et l'extrapolation.
  \item \textbf{Apprentissage fédéré} multi-sites : mutualiser l'apprentissage entre usines
        sans partager les données brutes (confidentialité), pour enrichir les signatures de
        défaut.
  \item Enrichissement du copilot RAG par la documentation machine du client (manuels,
        historique GMAO).
\end{itemize}

\subsection{Long terme --- solution industrielle complète}
\begin{itemize}[leftmargin=1.5em]
  \item Formalisation complète de l'architecture \textbf{OSA-CBM (ISO~13374)} et
        certification des processus de traçabilité.
  \item Connecteurs GMAO natifs (SAP PM, Maximo) au-delà du webhook générique.
  \item Application mobile pour les techniciens de terrain, export natif vers
        Grafana/Prometheus pour la supervision.
  \item \textbf{Place de marché de signatures de défaut} alimentée par la communauté, créant
        un effet de réseau autour du caractère ouvert de la plateforme.
\end{itemize}
